\section{Finite volume: a charge on vertex--face incidences}
\label{fv:section}

Count each corner once at its face and once at its vertex. We seek a
single corner charge that pays the face demands while respecting the
vertex allowances. The resulting charge is a dual certificate for the
necessary incidence problem. Searching for it in a specified family is
a second LP; its optimum is a property of that family, rather than a
geometric assumption.

\subsection{Vertex figures and the incidence allowance}
\label{fv:incidence}

Fix an integer $n\ge58$. A \emph{corner} is an incidence
$\kappa=(v,F)$ of a vertex $v$ and a two-dimensional face $F$.
The number of sides of $F$ is denoted by $s(F)$, and
\begin{equation}\label{fv:defect}
 D_5(P)=\sum_{\dim F=2}(5-s(F)).
\end{equation}
A triangle contributes two, a quadrilateral contributes one, and a
face with at least five sides contributes a nonpositive number.

We use the standard vertex-figure description of a finite-volume
Coxeter polytope. At a finite vertex the figure is a simplex. At an
ideal vertex a horospherical figure is a product
\begin{equation}\label{fv:link}
 L_v=\prod_{i=1}^k\Delta^{r_i},\qquad
 r_i\ge1,\qquad \sum_i r_i=n-1.
\end{equation}
Here $r_i$ is the dimension of a simplex factor, not the number of
nodes of its affine diagram. That diagram has $r_i+1$ nodes and Gram
rank $r_i$. These vertex-figure statements, including their applicability
to nonsimple ideal vertices, are retained geometric inputs from
\cite{BoundsArticle,BoundsSupplement}.

Put $r=n-1$. Write the $h$-polynomial of the figure as
\begin{equation}\label{fv:hpolynomial}
 h_v(t)=\sum_{j=0}^r h_{v,j}t^j,
 \qquad
 h_v(t)=\prod_i(1+t+\cdots+t^{r_i})
 \quad\hbox{at an ideal vertex}.
\end{equation}
At a finite vertex $h_{v,j}=1$ for every $0\le j\le r$.
Define
\begin{equation}\label{fv:allowance-general}
 V_v=\sum_{j=0}^r h_{v,j},\qquad
 G_v=\sum_{j=0}^r h_{v,j}\min(j,r-j),\qquad
 b_n(v)=\frac52G_v-\frac r2V_v.
\end{equation}
The quantity $V_v$ is the number of vertices of the vertex figure.
It equals $n$ at a finite vertex and $\prod_i(r_i+1)$ at an ideal
vertex. Every vertex of the figure has degree $r$, so its number of
edges is $E_v=rV_v/2$. These edges correspond to corners at $v$.

For a finite vertex, the sum in \eqref{fv:allowance-general} gives
\begin{equation}\label{fv:finite-allowance}
 \beta_n=b_n(v)=
 \begin{cases}
 n(n-6)/8,&n\text{ even},\\
 (n-1)(n-5)/8,&n\text{ odd}.
 \end{cases}
\end{equation}
At an ideal vertex we write
\begin{equation}\label{fv:ideal-allowance}
 K_n(v)=\frac{5\sum_jh_{v,j}\min(j,n-1-j)-(n-1)V_v}{2}.
\end{equation}
Thus $K_n(v)=b_n(v)$; the separate letter only distinguishes the
ideal allowance from the constant finite allowance $\beta_n$.

The incidence input needed below has the following precise form.

\begin{assumption}[Incidence excess]\label{fv:incidence-excess}
Let $P\subset\HH^n$ be a finite-volume Coxeter polytope, with the
vertex figures described above. The incidence theorem of
\cite{BoundsSupplement} gives
\begin{equation}\label{fv:incidence-excess-eq}
 D_5(P)\ge\sum_v b_n(v)+\Delta_n,
 \qquad \Delta_n\ge5(n-1)>0.
\end{equation}
\end{assumption}

This imports the height-index argument, including its strict excess,
as one geometric input. The LP need only prove
$D_5(P)\le\sum_vb_n(v)$; it need not produce a positive charge margin.

\subsection{The incidence feasibility problem and its charge certificate}
\label{fv:incidence-primal}

We first give the formulation without choosing particular corner
features. Fix a vector of rational features $\Psi(\kappa)\in\Q^d$.
The same feature vector must be used when a corner is viewed from its
face and when it is viewed from its vertex. A vertex type $V$ includes
its complete decorated vertex diagram or product-link data. A small-face
type $Q$ includes a triangle or quadrilateral together with its coherent
corner descriptions. Set
\begin{equation}\label{fv:feature-sums}
 A(V)=\sum_{\kappa\text{ at }V}\Psi(\kappa),\qquad
 A(Q)=\sum_{\kappa\text{ on }Q}\Psi(\kappa),\qquad
 d(Q)=5-s(Q).
\end{equation}
All domains used here are necessary domains: every actual vertex and
small face must have an allowed description. An allowed description
need not extend to an actual polytope.

Introduce nonnegative variables $x_V$ for vertex types, $y_Q$ for
small-face types, and $z_\kappa$ for allowed individual corners. The
last variables account for corners belonging to faces with at least
five sides. The incidence feasibility problem is
\begin{align}
 \sum_VA(V)x_V-\sum_QA(Q)y_Q-
       \sum_\kappa\Psi(\kappa)z_\kappa&=0,
       \label{fv:primal-balance}\\
 \sum_Qd(Q)y_Q-\sum_Vb_n(V)x_V&=1.
       \label{fv:primal-normalization}
\end{align}
If a domain is infinite, the variables may be read as finitely
supported nonnegative measures. Actual polytopes use only finitely
many incidences, which is all that the argument needs.

\begin{lemma}\label{fv:actual-lift}
An actual finite-volume Coxeter polytope gives a feasible point of
\eqref{fv:primal-balance}--\eqref{fv:primal-normalization} whenever
the allowed domains contain all its local descriptions.
\end{lemma}

\begin{proof}
Let $f_3$ and $f_4$ be the numbers of triangular and quadrilateral
two-faces, and put
\[
 L=2f_3+f_4-\sum_vb_n(v).
\]
Since the other two-faces contribute nonpositively to $D_5(P)$,
Input~\ref{fv:incidence-excess} gives
$L\ge\Delta_n>0$. Divide each actual vertex count, small-face count
and residual-corner count by $L$. Counting the same corners by vertices
and by faces gives \eqref{fv:primal-balance}. The definition of $L$
gives \eqref{fv:primal-normalization}.
\end{proof}

The residual variables in \eqref{fv:primal-balance} are important.
Their contribution has not disappeared merely because their faces have
nonpositive demand. Feature conservation still counts those corners.

\begin{proposition}[Charge certificate]\label{fv:charge-certificate}
Suppose a coefficient vector $\theta\in\R^d$ satisfies
\begin{align}
 b_n(V)-\theta\cdot A(V)&\ge0 &&\text{for every allowed vertex type},
    \label{fv:vertex-certificate}\\
 \theta\cdot A(Q)-d(Q)&\ge0 &&\text{for every allowed small face},
    \label{fv:face-certificate}\\
 \theta\cdot\Psi(\kappa)&\ge0 &&\text{for every allowed corner}.
    \label{fv:corner-certificate}
\end{align}
Then the incidence feasibility problem is infeasible.
\end{proposition}

\begin{proof}
Use Theorem~\ref{thm:homogeneous-certificate}, reversing the sign of
\eqref{fv:primal-balance}. The vertex, small-face and residual-corner
columns $(T,d)$ are respectively $(-A(V),-b_n(V))$,
$(A(Q),d(Q))$ and $(\Psi(\kappa),0)$; take $\lambda=\theta$.
The three displayed conditions are exactly its certificate inequalities.
\end{proof}

Equivalently, $w(\kappa)=\theta\cdot\Psi(\kappa)$ pays every triangle
and quadrilateral and stays within every vertex allowance. Faces with
at least five sides are already paid by nonnegative corner charges.
Summing gives $D_5(P)\le\sum_vb_n(v)$, contradicting incidence.
This is a direct weak-duality argument. It makes no use of
infinite-dimensional strong duality or a realization theorem for the
allowed local descriptions.

\subsection{The corner features used in the recovered certificate}
\label{fv:features}

We now specify the finite-dimensional family used in the actual
certificate. Its basic feature vector is
$\phi(\kappa)\in\Q^{26}$. A feature is a function of the decorated
corner, whereas its coefficient is an optimization variable.

At a finite vertex, the spherical diagram is a disjoint union of
connected finite diagrams. Its two marks are the two nodes determining
the corner. At an ideal vertex a corner corresponds to an edge of
the product \eqref{fv:link}: one simplex factor varies along an edge,
and a vertex is selected in every other factor. The varying affine
diagram has two marked nodes. Each other factor contributes the finite
diagram obtained by deleting the node selected by its vertex. These
finite residual diagrams are retained in the corner description.

We number $A_r$ along its path by $0,\ldots,r-1$. We number $D_r$
along the spine $0,\ldots,r-3$, with leaves $r-2,r-1$ attached to
$r-3$. The remaining finite and affine diagrams use the fixed
numberings of the rational tables in \cite{BoundsSupplement}.
One numbering is selected for a complete vertex diagram and used at
all its corners. A relabelling must carry the numbering with it; one
cannot choose a new favourable numbering separately at every corner.

Set $\phi_0=1$. Features $\phi_1,\phi_2,\phi_3$ are respectively
the prescribed ordinary pair weight for two marks in the same finite
component, the affine pair weight in the varying ideal factor, and
the ordinary weight for marks in different finite component copies;
each is zero outside its stated case. Their fixed rational tables,
including exceptional small ranks, are the intrinsic corner tables of
\cite{BoundsSupplement}. We use their final edition: in particular,
the ordinary $B_{13}$ endpoint entry is zero.

The next fifteen features are indicators. Each is one when the two
marks lie in a copy of the displayed component at the displayed
unordered pair, and when the final column holds. Otherwise it is zero.
A component mentioned as a witness is a connected component of the
complete corner diagram disjoint from the marked component. Put
$M_{15}=m_{A_{15}}+m_{D_{15}}$, the number of such finite components.
\begin{center}\small
\begin{tabular}{cllp{.30\textwidth}}
\toprule
$j$&Marked type&Marked pair&Witness condition\\\midrule
4&$A_{13}$&$\{0,12\}$&An $A_1$ occurs.\\
5&$A_{14}$&$\{0,13\}$&An $A_1$ occurs.\\
6,7&$A_{15}$&$\{0,13\},\{1,14\}$&An $A_1$ occurs.\\
8,9&$D_{15}$&$\{1,13\},\{1,14\}$&An $A_1$ occurs.\\
10,11&$A_{16}$&$\{1,12\},\{3,14\}$&An $A_1$ occurs.\\
12,13&$D_{16}$&$\{0,14\},\{0,15\}$&An $A_1$ occurs.\\
14,15&$D_{14}$&$\{0,12\},\{0,13\}$&An $A_1,A_2$ or $A_3$ occurs.\\
16,17&$D_{13}$&$\{0,11\},\{0,12\}$&$M_{15}\ge3$.\\
18&$D_{28}$&$\{0,12\}$&$M_{15}\ge2$.\\
\bottomrule
\end{tabular}
\end{center}
In a row with two feature indices, the two marked pairs correspond
in the displayed order.
The final seven indicators concern different finite component copies:
\begin{center}\small
\begin{tabular}{cp{.48\textwidth}p{.31\textwidth}}
\toprule
$j$&Marked components and allowed nodes&Witness condition\\\midrule
19&$A_4:\{0,3\}$ and $A_{12}:\{1,10\}$&An $A_1$ occurs.\\
20&$A_2:\{0,1\}$ and $A_4:\{0,3\}$&An $A_1$ and an $A_{10}$ occur.\\
21&$A_2:\{0,1\}$ and $D_5:\{0\}$&An $A_1$ and an $A_9$ occur.\\
22&$A_1:\{0\}$ and $D_5:\{0\}$&An $A_{11}$ occurs, or both an $A_9$ and an $A_2$ occur.\\
23&$A_1:\{0\}$ and $D_5:\{0\}$&An $A_9$ and an $A_2$ occur.\\
24&$A_{12}:\{0,11\}$ and $D_{16}:\{0\}$&$M_{15}\ge2$.\\
25&Two distinct $A_2$ copies, with either node in each&An $A_{10}$ and at least two $A_1$ copies occur.\\
\bottomrule
\end{tabular}
\end{center}
These are indicators of unordered marked pairs, not oriented-pair
counts. Features 22 and 23 deliberately overlap. Feature 25 vanishes
when the two marks lie in one $A_2$ copy. The raw features all lie in
$[0,1]$; this does not imply nonnegativity of an arbitrary linear
combination of them.

There is one additional ideal-corner indicator, denoted $\chi_B$.
Write $\Gamma[m]$ for an affine diagram with $m$ nodes and rank $m-1$.
The indicator is one precisely when the varying factor is $B[13]$,
the marked pair is $\{10,11\}$ or $\{10,12\}$, and the complete
residual finite diagram contains at least three components of types
$A_{15}$ or $D_{15}$. It is zero otherwise. Thus $B[13]$ denotes
standard type $\widetilde B_{12}$, not a finite rank-13 component.
The residual diagrams are formed by the other factors' one-node
deletions. One does not first delete the two marked nodes and count
newly separated pieces in the varying factor.

Let $\epsilon_n=[n=58]$. Use free vectors
$a^{58},a^{\rm tail}\in\R^{26}$ and parameters $\lambda,\rho\ge0$.
The searched charge is
\begin{equation}\label{fv:piecewise-charge}
 w(\kappa)=
 \bigl(\epsilon_na^{58}+(1-\epsilon_n)a^{\rm tail}\bigr)
       \cdot\phi(\kappa)+\epsilon_n\lambda\chi_B(\kappa).
\end{equation}
Its 53 conserved features are
\begin{equation}\label{fv:conserved-features}
 \Psi(\kappa)=
 (\epsilon_n\phi(\kappa),(1-\epsilon_n)\phi(\kappa),
                     \epsilon_n\chi_B(\kappa)).
\end{equation}
The margin $\rho$ is a separate certificate variable, giving 54
shared variables in total. It is not another corner feature.
The coefficient vectors are free: some entries in the successful
certificate are negative. Nonnegativity is imposed on the resulting
charge on admitted corners, rather than on every coefficient.

The finite table construction used in the checked endpoint network
has a declared ambient range through rank 58. No assertion about
larger dimensions is obtained by running that constructor outside
its range. The meaning and validity of \eqref{fv:piecewise-charge}
above the endpoint come from the separately retained all-rank face
theorems and analytic vertex estimates described below.

\subsection{The geometric face inputs and their linear conclusions}
\label{fv:face-inputs}

The small-face theorems determine which marked components and original
supporting walls can occur together, and which estimates can pay one
face. We retain their complete domains as the following LP interface.

\begin{assumption}[Retained local face theorems]\label{fv:face-assumption}
Use the complete corner, triangle and quadrilateral coverage of
\cite{BoundsSupplement}, together with the common-coefficient
transfers recorded in \cite{LPRepository}. Their coefficient conditions
form the declared local linear system $\mathcal L$. It contains the
corner and triangle rows on both dimension branches; the guarded
ideal and finite $D_{13}$ conditions; the fork and additional
rank-16 conditions; and the complete retained quadrilateral
coefficient conditions, including both branch-specific $A_{13}$ rows.
For any feasible shared coefficients, with the stated auxiliary
certificate variables, these theorems imply
\eqref{fv:face-certificate}--\eqref{fv:corner-certificate} for every
actual corner, triangle and quadrilateral in every $n\ge58$.
When a local transfer uses equalities forced by the endpoint vertex
budgets, its hypotheses include the complete network conditions of
Section~\ref{fv:finite-network}; those equalities are not presumed as
independent geometric axioms.
\end{assumption}

The word ``conditions'' in this input refers to restrictions
on searched charge coefficients. It does not assert that a polytope
has to satisfy numerical inequalities in variables $a^{58}$ and
$a^{\rm tail}$. Those variables describe a proof we are constructing.
Nor are the conditions claimed necessary for every possible valid
charge. For example, the retained $A_{13}$ argument uses
\begin{equation}\label{fv:A13-rows}
 a^{58}_0+a^{58}_4\ge1,
 \qquad a^{\rm tail}_0+a^{\rm tail}_4\ge1.
\end{equation}
These rows are part of this specified sufficient certificate family.

The coverage retains actual wall identities, ordered marks, complete
maps into the fixed numbered diagrams, residual choices at ideal links,
and every original wall used by an auxiliary construction. These are
the applicability conditions of the imported theorem; an inverse-Gram
calculation alone does not supply them. Face payments from different
alternatives are never added as though they belonged to disjoint corners.

Two examples show the actual linear form. For either branch vector
$z=a^{58}$ or $a^{\rm tail}$, the retained conditions include
\begin{align}
 4z_0+z_3/16+z_{22}+z_{23}&\ge1,
       \label{fv:example-face-P}\\
 4z_0+z_3/2+z_{21}&\ge1.
       \label{fv:example-face-M}
\end{align}
The geometric domains on which these rows pay the face are part
of their respective local theorems. The coefficients are not inferred
from the appearance of a matching type name elsewhere in a face.
The electronic local matrix gives every retained row and its source
association. Its geometric validity is the content of the cited local
results, not a consequence of the matrix having rational entries.

\subsection{Finite vertices: an exact vector-cost network}
\label{fv:finite-network}

At a finite vertex the marks range over all unordered pairs of
nodes in a spherical forest. A formal forest has nonnegative integer
component multiplicities $m_T$ satisfying
\begin{equation}\label{fv:forest-rank}
                  \sum_T r_Tm_T=58,
\end{equation}
where $r_T$ is the rank of type $T$. We use the finite collection of
weight types of the rational catalogue. A symbolic dihedral family
can be represented by one weight type only when its complete feature
table is the same throughout the stated parameter range and that
range is covered by the local theorem. A finite sample of labels is
not used in place of an unbounded family.

Let
\[
 \Phi(m)=\sum_{\kappa\text{ at the formal vertex}}\phi(\kappa).
\]
Every actual finite vertex supplies an admitted formal forest. The
converse is not required. We seek the family of inequalities
\begin{equation}\label{fv:finite-budget-endpoint}
                   a^{58}\cdot\Phi(m)+\rho\le377
                  \quad\hbox{for every }m\text{ satisfying
                  \eqref{fv:forest-rank}}.
\end{equation}
The right side is the incidence allowance
$\beta_{58}=58(58-6)/8$, rather than an assumed optimized vertex cost.

Here is the structure that makes this large family a network LP.
For one component $T$, let $S_T$ be the sum of its ordinary
same-component entries; for different component types $T,U$, let
$B_{TU}$ be the sum of their ordinary cross entries. Using $u_j$
for the $j$th coordinate vector of $\Q^{26}$ gives
\begin{equation}\label{fv:first-aggregates}
 \Phi_0=\binom{58}{2},\qquad
 \Phi_1=\sum_Tm_TS_T,\qquad
 \Phi_2=0,\qquad
 \Phi_3=\sum_{T<U}m_Tm_UB_{TU}.
\end{equation}
Features 4--18 sum to the marked-type multiplicity times their witness
indicator. The remaining coordinates follow by counting the unordered
mark pairs in the table: for example,
$\Phi_{19}=4m_{A_4}m_{A_{12}}[m_{A_1}>0]$ and
$\Phi_{25}=4\binom{m_{A_2}}2[m_{A_{10}}>0,m_{A_1}\ge2]$.
Brackets denote Boolean indicators; the factors four count node choices,
not orientations.

Fix the multiplicities $h_H=m_H$ of the types in
\begin{equation}\label{fv:high-set}
 \mathcal H=\{A_9,\ldots,A_{17},B_9,\ldots,B_{12},
                    D_9,\ldots,D_{16},D_{28}\}.
\end{equation}
Also fix $a=m_{A_1}$ and $b=m_{A_2}$; call the resulting tuple an
interface $I$. Let
\[
 M_h=58-\sum_Hr_Hh_H,\qquad R_I=M_h-a-2b.
\]
All nonzero ordinary cross terms have exactly one endpoint in
$\mathcal H$. After an interface is fixed, the remaining vector cost
is additive in the residual component counts, except for one
$A_3$-presence condition. More precisely, let $\mathcal R$ be the
remaining type alphabet after removing $\mathcal H,A_1,A_2$, and put
$\mathcal S=\mathcal R\setminus\{A_3\}$. Define
\[
 p_T^h=S_Tu_1+\Bigl(\sum_Hh_HB_{TH}\Bigr)u_3.
\]
For residual types other than $A_4,D_5$ use $p_{I,T}=p_T^h$.
The two exceptions have the complete vector costs
\begin{align}
 p_{I,A_4}={}&p_{A_4}^h+4h_{A_{12}}[a>0]u_{19}
                    +4b[a>0,h_{A_{10}}>0]u_{20},\notag\\
 p_{I,D_5}={}&p_{D_5}^h+2b[a>0,h_{A_9}>0]u_{21}\notag\\
 &+a[h_{A_{11}}>0\text{ or }(h_{A_9}>0,b>0)]u_{22}
                    +a[h_{A_9}>0,b>0]u_{23}.
                    \label{fv:residual-costs}
\end{align}
Let $K_I$ contain $\binom{58}{2}u_0$ and all terms of the feature
sum involving only the fixed counts $h,a,b$, including the fixed
same-component indicators. Put
\[
              \gamma_I=h_{D_{14}}[a=b=0](u_{14}+u_{15}).
\]
Then, with $k=m_{A_3}$,
\begin{equation}\label{fv:vector-cost-identity}
 \Phi(m)=K_I+kp_{I,A_3}+[k>0]\gamma_I
                         +\sum_{T\in\mathcal S}m_Tp_{I,T}.
\end{equation}
This identity follows directly from the feature definitions: the only
unfixed witness is $A_3$ in the $D_{14}$ indicators, the unfixed cross
terms are linear in $A_4,D_5$, and the others are fixed by $I$.

An entry for $(I,k)$ has vector cost
$K_I+kp_{I,A_3}+[k>0]\gamma_I$ and residual rank $R_I-3k$.
A transition removes a type $T\in\mathcal S$, reduces rank by
$r_T$, and contributes vector cost $p_{I,T}$. Only rank-feasible
interfaces, entries and reachable residual states occur. Rank strictly
decreases along every transition, so the graph is acyclic.

\begin{lemma}[Exact network formulation]\label{fv:network-exact}
Every admitted formal forest has a path whose vector cost is
$\Phi(m)$, and every path gives admitted counts of total rank 58.
Consequently the maximum path charge equals
$\max_m a^{58}\cdot\Phi(m)$ for every signed vector $a^{58}$.
\end{lemma}

\begin{proof}
The interface, entry and residual transitions select exactly the
nonnegative counts in \eqref{fv:forest-rank};
\eqref{fv:vector-cost-identity} gives the path cost. Different removal
orders can repeat a forest but cannot change its cost. A unit flow on
this acyclic graph is a convex combination of complete paths, so its
maximum linear cost is the maximum path cost.
\end{proof}

For a fixed high assignment, interfaces may share a residual network
only if their complete vectors $p_{I,A_4},p_{I,D_5}$ agree. The other
residual costs are already fixed by $h$. Such sharing retains every
entry separately. Equality of scalar costs at one previously selected
charge would not justify sharing while the charge coefficients vary.

For a network fragment $S$, write $v_{S,g,\ell}$ for free potentials
at reachable residual rank $\ell$ in sharing group $g$. Let $U_S$ be a
free fragment upper bound. The rows are
\begin{align}
 v_{S,g,0}&=0,\notag\\
 v_{S,g,\ell}-v_{S,g,\ell-r_T}
                  &\ge a^{58}\cdot p_{g,T},\notag\\
 U_S-v_{S,g(I),R_I-3k}
                  &\ge a^{58}\cdot(K_I+kp_{I,A_3}+[k>0]\gamma_I),
                  \label{fv:network-lp}\\
 U_S+\rho&\le377.\notag
\end{align}
All fragments use the same $a^{58},\rho$ and retain every residual
type; only potentials are private. Telescoping gives every path bound.
Conversely, maximum residual-completion costs give feasible potentials.
Eliminating them therefore gives exactly
\eqref{fv:finite-budget-endpoint}, even for signed coefficients.
Neither integral optimization variables nor nonnegative potentials
are needed. Completeness of the checked graph is a finite-data input;
the lemma explains why its encoding is an LP.

\subsection{Ideal vertices and the analytic tail}
\label{fv:ideal-block}

Let an ideal link have factor ranks
$r_1,\ldots,r_k$, total rank $R=n-1$ and $V=\prod_i(r_i+1)$ vertices.
For a connected affine factor $\Gamma$ of rank $r$, let $Q_\Gamma$
be the sum of its ordinary marked-pair weights and put
$q_\Gamma=Q_\Gamma/(r+1)$. Define $C=\sum_iq_{\Gamma_i}$.
The normalized ordinary charge for coefficient values $(a_0,a_2)$ is
\begin{equation}\label{fv:ideal-charge}
              \frac{C(v)}{V}=\frac R2a_0+Ca_2
\end{equation}
before adding the guarded indicator. This follows because a marked
pair in one factor occurs once for each choice of a vertex in every
other factor. Only raw features 0 and 2 occur at ideal corners.

Let $J$ be the sum of independent uniforms on
$\{0,\ldots,r_i\}$ and put $Y=2J-R$. The coefficient distribution
of $h_v$ is the distribution of $J$ multiplied by $V$. Symmetry gives
\begin{equation}\label{fv:normalized-ideal}
 \frac{K_n(v)}V=\frac{3R}{4}-\frac54\mathbb E|Y|,
 \qquad
 S_a=\frac{3R}{4}-\frac54\mathbb E|Y|-\frac R2a_0-Ca_2.
\end{equation}
The desired vertex inequality is $S_a\ge\rho$, so the raw ideal
margin is $V\rho$, not an unscaled copy of $\rho$.

\paragraph{The rank-57 block.}
For $n=58$, introduce free factor-cost variables $k_r$ for
$1\le r\le57$. For every connected affine family impose
\begin{equation}\label{fv:ideal-factor-links}
                 k_r\ge\frac r2a^{58}_0+q_\Gamma a^{58}_2.
\end{equation}
The complete family ranges are affine $A$ at ranks 1--57,
$C$ at 2--57, $B$ at 3--57, $D$ at 4--57, and the exceptional
$E_6,E_7,E_8,F_4,G_2$. There are 227 individual-family rows.
They all remain, even when $a^{58}_2$ is negative.

The ideal input supplies 354 retained symmetric payoff vectors, with
fixed ancestry and domination indices covering every factor addition.
Domination compares all central-block averages with those of one
retained vector; its index cannot vary with the coordinate. The indices
are data, while the payoff coordinates below are LP variables. The
underlying operation on an array is uniform convolution minus its
factor cost:
\[
 f'_j=\frac1{r+1}\sum_{\ell=0}^r f_{j+\ell}-k_r.
\]
Its initial array is $g_j=(5\min(j,57-j)-57)/2$; its final scalar
is the normalized allowance minus the factor costs. At accumulated
rank $R$, put $L=57-R$ and let
$B_{R,i,a}$ be the payoff average on the
central interval $a,\ldots,L-a$, for
$0\le a\le\lfloor L/2\rfloor$. Let $I_R$ be the retained index
set at level $R$. If a rank-$r$ factor is added to parent
$j\in I_{R-r}$, put $A=L-2a+1$ and
\begin{align}
 \omega_{L,a,r,\ell}&=\frac{A+r-2\ell}{A(r+1)},
       &&0\le\ell<\min(A,r+1),\notag\\
 T_{R,i,a}^{r,j}&=B_{R,i,a}+k_r-
       \sum_{\ell=0}^{\min(A,r+1)-1}
           \omega_{L,a,r,\ell}B_{R-r,j,a+\ell}.
       \label{fv:ideal-average-transform}
\end{align}
These positive coefficients sum to one and express the central average
of the convolution as an average of nested central intervals. The rows
are:
\begin{enumerate}[label=\textup{(\roman*)},nosep]
\item the 29 initial equalities
\[
 B_{0,0,a}=\frac1{58-2a}\sum_{j=a}^{57-a}
                  \frac{5\min(j,57-j)-57}{2},
 \qquad0\le a\le28;
\]
\item $T_{R,i,a}^{r,j}=0$ for the fixed ancestry $(r,j)$ of
 each retained $i\in I_R$;
\item $T_{R,D(R,r,j),a}^{r,j}\le0$ for every candidate, using
 its fixed retained dominator $D(R,r,j)$;
\item $\tau\le B_{57,i,0}$ for every final retained index, and
 $\rho\le\tau$, with $\tau$ free.
\end{enumerate}
The fixed index sets and maps are specified in the retained electronic
ideal data of \cite{LPRepository}. With the family links this block has
103,265 rows and 5,754 private variables. Positive convolution propagates
the retained lower bounds through every factor addition, so every
rank-57 profile has surplus at least $\tau\ge\rho$. These fixed
domination choices define a sufficient certificate family, rather than
all possible ideal-vertex laws.

\paragraph{The guarded ideal feature.}
The additive block just described remains in the LP. The additional
feature $\chi_B$ is handled separately; replacing its effect by an
unconditional change of the rank-12 factor cost would be incorrect.
If $\chi_B$ occurs, the varying $B[13]$ factor has rank 12 and
the residual finite rank is 45. Three rank-15 components exhaust that
rank. The complete one-node-deletion classification leaves exactly
\begin{align}
 \mathrm{I}:&\quad B[13]\times\Gamma_1[16]
                   \times\Gamma_2[16]\times\Gamma_3[16],\notag\\
 \mathrm{II}:&\quad B[13]\times D[31]\times\Gamma[16],
 \qquad\Gamma_i,\Gamma\in\{A,B,D\}.
 \label{fv:guarded-links}
\end{align}
Let $a_\chi=V^{-1}\sum_\kappa\chi_B(\kappa)$. The additional rows are
\begin{equation}\label{fv:guarded-rows}
  \frac{57}{2}a^{58}_0+C a^{58}_2+a_\chi\lambda+\rho\le K.
\end{equation}
All thirteen rows have the following compact specification. Put
\[
 (c_A,c_B,c_D)=(5,25/4,411/64),\qquad
 (d_A,d_B,d_D)=(1,1/16,1/4).
\]
For every unordered triple $\Gamma_1,\Gamma_2,\Gamma_3\in\{A,B,D\}$
in type I, and every $\Gamma\in\{A,B,D\}$ in type II, use
\begin{equation}\label{fv:guarded-data}
\begin{array}{c|ccc}
 &C&a_\chi&K\\\hline
 \mathrm{I}&11/2+\sum_{i=1}^3c_{\Gamma_i}
       &(2/13)\prod_{i=1}^3d_{\Gamma_i}&331643/13312\\
 \mathrm{II}&753/62+c_\Gamma&(2/403)d_\Gamma&16549/806.
\end{array}
\end{equation}
Here $d_\Gamma$ is the proportion of one-node deletions giving the
required residual component; the deletion counts are $16,1,4$.
The values of $K$ are \eqref{fv:normalized-ideal} for ranks
$(12,15,15,15)$ and $(12,30,15)$. Thus the specification retains
deletion multiplicities, not just the affine factor types.

\paragraph{The residual dimension-59 sector.}
Append all 63 connected-family links of ranks at most 16 for the
tail coefficients,
\begin{equation}\label{fv:tail-factor-links}
             k_r\ge\frac r2a^{\rm tail}_0
                                      +q_\Gamma a^{\rm tail}_2,
\end{equation}
using the \emph{same} costs $k_r$ as above. Put $k_0=0$ and introduce
$\mu\ge0$. The fifteen descent rows are
\begin{equation}\label{fv:descent}
 k_r-k_{r-1}+\mu\le h_r,\qquad 1\le r\le15,
\end{equation}
where
\[
 (h_1,\ldots,h_{15})=
 \left(\frac34,\frac{47}{66},\frac{89}{132},\frac7{11},
 \frac{79}{132},\frac{37}{66},\frac{23}{44},
 \frac{2213}{4374},\frac{59}{132},\frac9{22},
 \frac{67}{156},\frac{71}{168},\frac5{12},
 \frac{79}{192},\frac{83}{204}\right).
\]
The retained descent theorem states that a rank-58 product with
largest factor rank between 10 and 16 has a factor whose rank can be
reduced by one so that its envelope surplus increases on restoring
that factor by at least $\mu$. The reduced product has rank 57 and
is covered by the full preceding block. This proves the required
dimension-59 ideal budget in that sector. Independent copies of the
costs in the two blocks would not justify this argument.

\paragraph{The unbounded ideal tail.}
For every remaining dimension-59 ideal link, and every ideal link
in dimension at least 60, the retained analytic theorem gives surplus
at least $\varepsilon=1/5307$ for baseline coefficients $(0,1)$.
For an arbitrary tail law a sufficient interval test reduces exactly
to four affine rows:
\begin{align}
 a^{\rm tail}_0&\le0,&
 a^{\rm tail}_0+a^{\rm tail}_2&\le1,\notag\\
 \rho+29a^{\rm tail}_0&\le\varepsilon,&
 \rho+29a^{\rm tail}_0+29a^{\rm tail}_2&\le29+\varepsilon.
 \label{fv:analytic-ideal-rows}
\end{align}
To see the reduction, write $t=1/n$,
$u_0=a^{\rm tail}_0$, $u_1=a^{\rm tail}_0+a^{\rm tail}_2-1$ and
\[
 P_i(t)=(\varepsilon-\rho)t-\frac{1-t}{2}u_i.
\]
Each $P_i$ is affine, and \eqref{fv:analytic-ideal-rows} is exactly
nonnegativity at the two endpoints of $[0,1/59]$. The analytic input
also gives $0\le C\le R/2$. If $w=2C/R$, then
\[
 t(S_a-\rho)=t(S_{(0,1)}-\varepsilon)
                         +(1-w)P_0(t)+wP_1(t)\ge0.
\]
Thus these four rows prove the desired budget on the entire stated
tail. Their validity uses the analytic surplus theorem; it is not
inferred from checking finitely many dimensions.

\paragraph{Finite vertices above the endpoint.}
Let $q\in\Q^{26}$ be the vector displayed in
\eqref{fv:successful-law} below. Introduce $0\le\xi\le1$ and impose
\begin{equation}\label{fv:finite-tail-rows}
 a^{\rm tail}_j\le q_j+2\xi[\,j\in\{16,17\}\,]
       \quad(j\ne2),\qquad \rho\le\frac14.
\end{equation}
The retained analytic finite-vertex estimate asserts that, for every
formal finite forest of every rank $n\ge59$,
\[
 \beta_n-\bigl(q+2\xi(u_{16}+u_{17})\bigr)\cdot\Phi_n
                                                       \ge\frac14.
\]
Since finite raw features are nonnegative and feature 2 is zero,
\eqref{fv:finite-tail-rows} proves its vertex budget. This is a
spectator- and rank-uniform analytic input, separate from the rank-58
catalogue. The coordinatewise comparison is a sufficient restriction
on the searched tail law, not a necessary description of all valid
finite-vertex charges.

\subsection{Optimization of the finite-vertex budget}
\label{fv:377}

The finite-vertex problem asks for the smallest allowance required
by a charge family. Once a finite formal forest domain $\mathcal T$ and a
local coefficient family $\mathcal C$ have been specified, write
\begin{equation}\label{fv:budget-template}
 U_{\mathcal T}(a)=\max_{F\in\mathcal T}a\cdot\Phi(F),
 \qquad \min_{a\in\mathcal C}U_{\mathcal T}(a).
\end{equation}
Here $a$ is the coefficient vector used at the finite vertex,
$\Phi(F)$ is its aggregate feature vector, and $\mathcal C$ is the
projection of the specified joint certificate family onto $a$.
The auxiliary variables and their conditions remain part of that
joint family.
The relevant geometric allowance is supplied independently by
incidence. A complete feasible certificate closes the finite-vertex
part when its required allowance is no greater than that geometric
allowance. Its face and ideal-vertex obligations, including all
dependencies of their local proofs, must hold at the same shared
coefficients.

Here $\mathcal T$ is the complete admitted rank-58 domain, and the
independently supplied allowance is $\beta_{58}=377$. The following
specialization gives both a feasible law and an exact dual lower bound
for its declared coefficient family. It makes no extension of the
endpoint catalogue to other ranks.

The following rational point uses the coordinate order specified in
Section~\ref{fv:features}:
\begin{align}
 (q_0,\ldots,q_{12})={}&
 (0,1,1,1,1,\tfrac12,\tfrac14,\tfrac14,\tfrac14,\tfrac14,
       \tfrac58,\tfrac58,\tfrac12),\notag\\
 (q_{13},\ldots,q_{25})={}&
 (\tfrac12,\tfrac12,\tfrac12,-\tfrac12,-\tfrac12,\tfrac12,
       \tfrac18,\tfrac14,\tfrac12,\tfrac38,\tfrac58,\tfrac12,
       \tfrac14),\notag\\
 a^{58}={}&q,\qquad
 a^{\rm tail}=q+\tfrac12(u_{16}+u_{17}),\notag\\
 \lambda={}&\tfrac12,\qquad\rho=0.
 \label{fv:successful-law}
\end{align}
In particular the two $D_{13}$ indicator coefficients are negative at
the endpoint and zero on the tail. This does not contradict the
nonnegativity of the resulting corner charge, which is checked on
the admitted decorated-corner domain.

The private variables can also be specified exactly. In the finite-tail
and descent blocks take $\xi=1/4$ and $\mu=13/2187$. Take
$k_r=\max_{\operatorname{rank}\Gamma=r}q_\Gamma$, over the complete
connected-family inventory, and use the retained ideal payoff vectors
and their ancestry equalities to assign every $B_{R,i,a}$, with
$\tau=0$. This is an assignment to the complete ideal certificate,
not a selection of a numerically best product profile. The actual
complete verification also checks every finite-network potential at
these same shared coefficients.

For example, the face rows
\eqref{fv:example-face-P}--\eqref{fv:example-face-M} take values
$17/16$ and $1$ on both branches. The guarded ideal rows have slack
at least $939/13312$ in type I and $50623/25792$ in type II. These
examples illustrate the assignment; complete exact verification checks
every local and network row at the same shared coefficients.

There is a useful optimization statement in addition to feasibility.
Let $\mathcal A$ be the projection onto
$(a^{58},a^{\rm tail},\lambda)$ of the complete local coefficient
system, including \eqref{fv:A13-rows}, with $\rho=0$. Replace the
last network row by $U_S\le U$, for one free common variable $U$,
and minimize $U$. Eliminating the potentials gives
\begin{equation}\label{fv:minimax}
 \min_{(a^{58},a^{\rm tail},\lambda)\in\mathcal A}
                 \max_F a^{58}\cdot\Phi(F),
\end{equation}
where $F$ ranges over the complete admitted rank-58 formal-forest
domain. The allowance 377 is no longer imposed as an upper bound
in this diagnostic optimization.

\begin{proposition}[Exact restricted minimax]\label{fv:exact-minimax}
The optimum in \eqref{fv:minimax} is 377.
\end{proposition}

\begin{proof}
The complete feasible assignment \eqref{fv:successful-law}, including
its private local variables and all network potentials, gives an upper
bound of 377. For the converse, write $a_0,a_1$ for coordinates 0 and
1 of $a^{58}$. The retained corner and triangle rows give
$a_0\ge0$ and $3a_0+2a_1\ge2$.
The formal forest
\[
                    F_*=2D_{17}+D_{24}
\]
has total rank 58. The ordinary table gives
$S_{D_{17}}=227/2$ and
$S_{D_{24}}=8\cdot24-32-1-4-3-2=150$. All ordinary cross
and contextual features vanish on this forest, so its aggregate
feature vector is
$\Phi(F_*)=1653u_0+377u_1$. Thus one path bound gives
$U\ge1653a_0+377a_1$. The exact nonnegative combination
\begin{equation}\label{fv:minimax-identity}
 U-377=(U-1653a_0-377a_1)
 +\frac{377}{2}(3a_0+2a_1-2)+\frac{2175}{2}a_0
\end{equation}
proves the lower bound: the first term is the sum of this forest's
network rows and its link to $U$, and the other terms are multiples
of the two local rows.
\end{proof}

The allowance 377 comes from incidence; the optimum 377 comes from
feasibility and \eqref{fv:minimax-identity}. This is optimality within
the declared law family and formal domain. The forest $F_*$ need not
be geometrically realized, and richer valid law families may have
smaller optima.

In the original capped system,
$\max_Fa^{58}\cdot\Phi(F)+\rho\le377$ and $\rho\ge0$.
The same inequalities force $a_0=0$, $a_1=1$ and $\rho=0$.
Thus maximizing this margin and requiring it to be positive would
miss the successful certificate. Incidence already has the strict
excess needed for the proof.

\subsection{All dimensions and the recovered bound}
\label{fv:all-dimensions}

For clarity, the separate coverage obligations are collected below.
They use one dimension-dependent charge, not independent face and
vertex laws.
\begin{center}\small
\begin{tabular}{p{.29\textwidth}p{.61\textwidth}}
\toprule
Domain&Retained certificate argument\\\midrule
Finite vertices, $n=58$&The complete vector-cost network and all its
 entry and transition inequalities.\\
Finite vertices, $n\ge59$&The analytic finite-vertex estimate and
 coordinatewise domination \eqref{fv:finite-tail-rows}.\\
Ideal vertices, $n=58$, $\chi_B$ aggregate zero&The complete rank-57
 ideal block and connected-family links.\\
Ideal vertices, $n=58$, positive $\chi_B$ aggregate&The complete
 deletion classification \eqref{fv:guarded-links} and all thirteen
 rows \eqref{fv:guarded-rows}.\\
Ideal vertices, $n=59$, largest factor rank 10--16&The same costs
 and the complete rank-descent inequalities.\\
Other ideal vertices at 59, and all $n\ge60$&The analytic ideal
 surplus theorem and the four rows \eqref{fv:analytic-ideal-rows}.\\
All actual corners and small faces, $n\ge58$&The complete local
 geometric and contact coverage of Input~\ref{fv:face-assumption}.\\
\bottomrule
\end{tabular}
\end{center}

\begin{theorem}[Finite-volume recovery]\label{fv:recovery}
The retained local geometric, classification, incidence and exact
finite-data theorems, together with the verified common-coefficient
certificate above, exclude every finite-volume hyperbolic Coxeter
polytope in every integer dimension $n\ge58$.
\end{theorem}

\begin{proof}
Suppose such a polytope exists. The last row of the table gives
nonnegative corner charges and the required payments on its triangles
and quadrilaterals. The preceding rows give
\[
 \sum_{\kappa\text{ at }v}w(\kappa)\le b_n(v)
\]
at every finite or ideal vertex. Apply the charge certificate and the
actual lift of Lemma~\ref{fv:actual-lift}.
\end{proof}

